\documentclass[11pt]{article}
\usepackage[utf8]{inputenc}
\usepackage{amsmath,amssymb,amsthm}
\input{glyphtounicode-cmex}
\pdfgentounicode=1
\usepackage[margin=1in]{geometry}
\usepackage{microtype,enumitem,booktabs}
\usepackage{color}
\definecolor{linkblue}{rgb}{0.05,0.15,0.35}
\PassOptionsToPackage{hyphens}{url}
\usepackage[colorlinks=true,linkcolor=linkblue,citecolor=linkblue,urlcolor=linkblue]{hyperref}
\usepackage[nameinlink,noabbrev,capitalise]{cleveref}
\hypersetup{pdftitle={A Cyclic Polytabloid Proof of Saxl's Conjecture},pdfauthor={OpenAI}}
\numberwithin{equation}{section}
\newtheorem{theorem}{Theorem}[section]
\newtheorem{lemma}[theorem]{Lemma}
\newtheorem{proposition}[theorem]{Proposition}
\newtheorem{corollary}[theorem]{Corollary}
\theoremstyle{definition}
\newtheorem{definition}[theorem]{Definition}
\newtheorem{example}[theorem]{Example}
\theoremstyle{remark}
\newtheorem{remark}[theorem]{Remark}
\AddToHook{env/theorem/begin}{\crefalias{theorem}{theorem}}
\AddToHook{env/lemma/begin}{\crefalias{theorem}{lemma}}
\AddToHook{env/proposition/begin}{\crefalias{theorem}{proposition}}
\AddToHook{env/corollary/begin}{\crefalias{theorem}{corollary}}
\AddToHook{env/definition/begin}{\crefalias{theorem}{definition}}
\AddToHook{env/example/begin}{\crefalias{theorem}{example}}
\AddToHook{env/remark/begin}{\crefalias{theorem}{remark}}
\newcommand{\C}{\mathbb C}
\DeclareMathOperator{\Span}{span}
\DeclareMathOperator{\Ind}{Ind}
\DeclareMathOperator{\Res}{Res}
\DeclareMathOperator{\Hom}{Hom}
\DeclareMathOperator{\sgn}{sgn}
\DeclareMathOperator{\id}{id}
\DeclareMathOperator{\GL}{GL}
\DeclareMathOperator{\Mat}{Mat}
\DeclareMathOperator{\adj}{adj}
\DeclareMathOperator{\tr}{tr}
\DeclareMathOperator{\Sym}{Sym}
\DeclareMathOperator{\Alt}{Alt}
\DeclareMathOperator{\supp}{supp}
\DeclareMathOperator{\diag}{diag}
\DeclareMathOperator{\len}{len}
\setlist[enumerate]{itemsep=3pt,topsep=5pt}
\setlist[itemize]{itemsep=3pt,topsep=5pt}
\setcounter{tocdepth}{2}
\title{A Cyclic Polytabloid Proof of Saxl's Conjecture}
\author{OpenAI}
\date{September 24, 2026}
\begin{document}
\maketitle
\begin{abstract}
For every staircase partition, we prove that the tensor square of the
corresponding irreducible complex representation of the symmetric group
contains every irreducible representation of that group. This proves
Saxl's conjecture.
\end{abstract}
\section{Introduction}
\label{sec:introduction}

For a partition $\alpha\vdash n$, let $S^\alpha$ denote the corresponding
irreducible complex representation of the symmetric group $S_n$.
The Kronecker coefficient
\[
g(\alpha,\beta,\lambda)
=\dim\Hom_{S_n}(S^\lambda,S^\alpha\otimes S^\beta)
\]
is the multiplicity of $S^\lambda$ in the tensor product, with $S_n$
acting diagonally. Put
\begin{equation}\label{intro:staircase}
N_m=\frac{m(m+1)}2,\qquad \rho_m=(m,m-1,\ldots,1).
\end{equation}
We prove the following statement.

\begin{theorem}[Saxl's conjecture]\label{thm:saxl}
For every integer $m\ge1$ and every partition $\lambda\vdash N_m$,
\[
g(\rho_m,\rho_m,\lambda)>0.
\]
Equivalently, $S^{\rho_m}\otimes S^{\rho_m}$ contains every irreducible
complex representation of $S_{N_m}$.
\end{theorem}

The assertion is uniform in both the staircase size and the constituent.
The proof establishes it inside one explicitly generated submodule of
the tensor square, a stronger conclusion that retains the vector needed
at each induction step.

The companion \cite[Theorem~1.1]{universaltensorsquares2026} uses the
stronger cyclic statement of \Cref{thm:cyclic-saxl}, together with
additional arguments, to construct, for every positive integer
$n\notin\{2,4,9\}$, an irreducible complex representation of $S_n$ whose
tensor square contains every irreducible representation of $S_n$.

\subsection{History and significance}

Saxl's conjecture belongs to a broader tensor-square covering problem:
can a single irreducible representation have every irreducible as a
constituent of its square? A motivating precedent comes from finite
simple groups of Lie type. Heide, Saxl, Tiep, and Zalesski proved that
the Steinberg square has this property, except for the groups
$\mathrm{PSU}_d(q)$ with $d\ge3$ and $\gcd(d,2(q+1))=1$
\cite[Theorem~1.2]{heidestz2013}.
Saxl proposed the staircase assertion in a UCLA Combinatorics Seminar on
20 March 2012, as recorded by Pak, Panova, and Vallejo
\cite[Conjecture~1.2 and Footnote~2]{pakpanovavallejo2016}.
Their work places the staircase square among tensor-square covering
questions for symmetric groups and proves several families of positive
Kronecker coefficients; in particular, their Theorem~1.4 covers hooks
and two-row partitions for sufficiently large staircase size.
Ikenmeyer proved positivity when $\lambda$ and $\rho_m$ are comparable
in the dominance order \cite[Theorem~2.1]{ikenmeyer2015}.
We use his triangular arrangement of two alternating tensor layers
\cite[Section~4]{ikenmeyer2015}; his argument also yields all hook constituents
\cite[Corollary~6.1]{ikenmeyer2015}.

Character-theoretic and modular methods have supplied further substantial
families. The character-value criterion of Pak, Panova, and Vallejo
\cite{pakpanovavallejo2016} detects constituents using the principal-hook
cycle type of a self-conjugate partition. Bessenrodt developed related
critical-class and spin-character methods, including the double-hook
case \cite{bessenrodt2018}. Li proved the triple-hook case using
computer-assisted finite reductions \cite{li2021}.
Here double and triple hooks are diagrams whose largest square has side
two and three, respectively.
Bessenrodt, Bowman, and Sutton connected Kronecker positivity to
2-modular decomposition numbers; their results include all
2-height-zero constituents, hence all odd-degree constituents, and
framed staircases \cite[Theorems~5.2 and~5.5]{bessenrodtbowmansutton2021}.

Another line of work addresses the tensor exponent and asymptotic
coverage. Luo and Sellke used the semigroup property to prove that a
random partition occurs in the staircase square with probability tending
to one, for both the uniform and Plancherel measures
\cite[Theorems~1.5 and~1.6]{luosellke2017}. They also found, in every
sufficiently large degree, an irreducible representation whose fourth
tensor power contains every irreducible
\cite[Theorem~1.4]{luosellke2017}. Harman and Ryba proved full coverage
for the tensor cube of every staircase, with combinatorial and modular
proofs \cite[Theorem~1.1]{harmanryba2023}.
Asymptotic coverage does not exclude exceptional constituents, while a
cube or fourth power uses additional tensor factors; neither gives the square
assertion in \Cref{thm:saxl}.

Recent developments include semigroup constructions of further
constituents and generalized-block constraints on support
\cite{ebrahimi2025}, and a full
staircase-square theorem for unipotent characters of $\GL_n(q)$
\cite[Theorem~1.1]{letelliernam2025}. The latter is an analogue in a
different character theory: ordinary Kronecker positivity implies the
corresponding unipotent positivity, but the converse need not hold.

The band calculation below also has a precise predecessor at the level
of ambient representations. Brown, van Willigenburg, and Zabrocki proved
that
$S^{(m,m-1)}\otimes S^{(m,m-1)}$ is the multiplicity-free sum of all
irreducibles of $S_{2m-1}$ indexed by partitions with at most four rows
\cite[Corollary~4.1]{brownwilligenburgzabrocki2010}.
Our induction needs each of these constituents in the cyclic module of
one prescribed band vector, rather than merely somewhere in that
ambient square. Explicit nonzero pairings supply this additional
generator-specific statement. Throughout, the objective is support,
or positivity, rather than a formula for Kronecker multiplicities.

\subsection{The construction and its reusable steps}

Place $N_m$ tensor positions in the triangle
$B_m=\{(i,j):i,j\ge1,\ i+j\le m+1\}$.
Let $R_m$ and $C_m$ be its row and column set partitions, ordered within
each row by increasing $j$ and within each column by increasing $i$.
Use the ordered basis $e_1,\ldots,e_m$ of $\C^m$. On a block of length
$r$, alternate the initial letters $e_1,\ldots,e_r$, taking the signed
sum over all permutations. Taking the product over row blocks gives
$v_{R_m}$; doing the same over column blocks gives $v_{C_m}$.
These are polytabloids of staircase type. Instead of beginning with an
arbitrary vector in the tensor square, we work throughout with
\[
W_m=\C[S_{B_m}](v_{R_m}\otimes v_{C_m}),
\]
where $S_{B_m}$ permutes the positions in both layers simultaneously.
\Cref{thm:cyclic-saxl} asserts that this particular cyclic module already
contains every irreducible. We first prove that it contains every
$S^\lambda$ for which $\lambda$ is dominated by $\rho_m$: each initial
sum of the parts of $\lambda$ is at most the corresponding staircase
sum. A sign twist and a permutation-module quotient give this starting
case (\Cref{prop:dominated}). The following three steps reach the
remaining partitions.

\begin{enumerate}
\item \textbf{A full induced quotient from a coordinate projection.}
Separate a smaller triangle $B_{m-s}$ from a band of width $s=1$ or $2$.
Call letters $1,\ldots,s$ low and letters $s+1,\ldots,m$ high.
Each position carries one letter from each tensor layer.
Project onto words whose two letters at each position are either both
high or both low. The staircase row and column counts force the high
positions of the projected generator to be exactly the smaller triangle.
The generator then factors into the smaller cyclic generator and a band
generator, whose cyclic module is denoted by $U_{m,s}$.
Translates of the high-position set occupy disjoint
coordinate sectors, so the quotient is the full induced module
\[
\Ind_{S_{B_{m-s}}\times S_{B_m\setminus B_{m-s}}}^{S_{B_m}}
       (W_{m-s}\boxtimes U_{m,s}).
\]
\Cref{lem:sector-induction} isolates the general reason that sector
separation gives an induced-module isomorphism, including an explicit
inverse. \Cref{prop:band-quotient} gives the resulting quotient.

\item \textbf{Constituents detected in the specified band generator.}
A width-two band is an odd path of alternating pairs with fixed endpoint
letters. At each position, a covector on the pair-letter space is
represented by a two-by-two matrix. Pairing along the path evaluates a word of
matrices and adjugates. To test a partition with at most four rows,
place its tableau columns on consecutive path segments and alternate
the covectors within each column. Each segment has length at most four.
A four-element basis of $\Mat_2(\C)$ makes every such alternated segment
matrix invertible.
A simultaneous conjugation of the test basis makes the endpoint entry
of the product nonzero. The resulting dual-polytabloid pairing proves
that every shape with at most four rows occurs in $U_{m,2}$ itself
(\Cref{prop:two-band}).

\item \textbf{A horizontal-strip reduction with four steps.}
If a shape outside this dominance range has a first row of length at
least $m$, one horizontal $m$-strip suffices. Otherwise its first four
rows contain more than $2m-1$ boxes. Successive sweeps of the bottom
boxes of columns then remove exactly $2m-1$ boxes in at most four
horizontal strips (\Cref{comb:strip-reduction}). In this second case,
let $\nu$ be the remaining partition. Iterated Pieri supplies a shape
$\eta\vdash2m-1$ with at most four rows such that $S^\lambda$ occurs
in the representation induced from $S^\nu\boxtimes S^\eta$.
The band calculation puts $S^\eta$ in $U_{m,2}$, while induction puts
$S^\nu$ in $W_{m-2}$. The full induced quotient then gives the target
constituent in $W_m$.
\end{enumerate}

The smaller factor in the quotient is the actual cyclic module $W_{m-s}$,
and the band factor is generated by the actual projected vector.
This is why both parts of the induction track generators instead of
using ambient tensor-product support alone. Only one copy of a suitable
intermediate irreducible is needed.

\Cref{sec:specht-tools} fixes the tensor conventions and proves the
needed form of Young's rule from Pieri. \Cref{sec:dominance} defines
$W_m$ and proves the dominated case. The quotient and band calculations
occupy \Cref{sec:band-quotient,sec:path-band};
\Cref{sec:completion} proves the strip reduction and both main theorems.

\section{Specht modules and horizontal strips}
\label{sec:specht-tools}

We fix the tensor conventions used to identify alternating block vectors
with Specht modules. The other input needed by the induction is the
Pieri rule: it transfers constituent support across the removal of
horizontal strips. We record its dominance consequences with proofs.

All representations and tensor products are over $\C$ and are
finite-dimensional.  A partition $\lambda\vdash n$ is padded with zero
parts when necessary; its length is $\len(\lambda)$ and its conjugate is
$\lambda^t$, where $\lambda^t_j=\#\{i:\lambda_i\ge j\}$.  We identify
$\lambda$ with its diagram $\{(i,j):1\le j\le\lambda_i\}$.
For partitions of the same size, write $\lambda\unrhd\mu$ if
$\sum_{i=1}^k\lambda_i\ge\sum_{i=1}^k\mu_i$ for every $k\ge1$.
Conjugation reverses this order:
\begin{equation}\label{tools:conjugate-dominance}
 \lambda\unrhd\mu\quad\Longleftrightarrow\quad
 \mu^t\unrhd\lambda^t.
\end{equation}
Indeed, for every integer $q\ge1$,
\[
 \sum_i\max(\lambda_i-q,0)
 =\max_{k\ge0}\left(\sum_{i=1}^k\lambda_i-kq\right),
 \qquad
 \sum_{j=1}^q\lambda^t_j
 =n-\sum_i\max(\lambda_i-q,0).
\]
Dominance increases the first expression and therefore decreases the
second.  Applying the resulting implication to the conjugate partitions
proves the converse.

\subsection{Position tensors and polytabloids}

For a finite set $B$, write $S_B$ for its symmetric group. For
$\lambda\vdash |B|$, write $S_B^\lambda$ for the irreducible $S_B$-module
of type $\lambda$; omit the subscript when the position set is understood.
If $E$ has
ordered basis $e_1,\ldots,e_d$, the tensor space $E^{\otimes B}$ has word
basis $e_a=\bigotimes_{b\in B}e_{a(b)}$, indexed by maps
$a:B\to\{1,\ldots,d\}$.  The position action is
$g e_a=e_{a\circ g^{-1}}$.  All regroupings of tensor factors are ordinary
tensor identifications and introduce no signs.

Let $L$ be a set partition of $B$ whose blocks have size at most $d$.
For each block $A=(b_1,\ldots,b_r)$ choose an ordering and set
\begin{equation}\label{tools:block-tensor}
 v_A=\sum_{\pi\in S_r}\sgn(\pi)
       \bigotimes_{i=1}^r e_{\pi(i)},\qquad
 v_L=\bigotimes_{A\in L}v_A,
\end{equation}
where the $i$th factor of $v_A$ occupies position $b_i$.
Changing a block ordering changes $v_L$ only by a sign.  Its nonzero
word coefficients are $\pm1$: on each block of size $r$, the letters are
a bijection with $\{1,\ldots,r\}$.

\begin{lemma}[Alternating blocks as polytabloids]\label{lem:polytabloid}
Define $\theta_i=\#\{A\in L:|A|\ge i\}$, so that the block sizes are
the column lengths of $\theta$.  Then
$\C[S_B]v_L\cong S^\theta$.
The same assertion holds for any ordered basis of a dual vector space.
\end{lemma}

\begin{proof}
Arrange the blocks as columns of a tableau $t$ of shape $\theta$, in
nonincreasing order of their sizes, and place each ordered block from top
to bottom.  Tableau entries are the distinct elements of $B$.
A row tabloid $\{t'\}$ records the set of entries in each row.
Let $M^\theta$ be the permutation module spanned by these row tabloids.
Send $\{t'\}$ to
the word assigning letter $i$ to every entry of row $i$.
This is an $S_B$-equivariant bijection from the tabloid basis of
$M^\theta$ to the word basis of content $\theta$.
If $C_t$ is the subgroup permuting entries within each column of $t$,
the polytabloid
\[
 e_t=\sum_{c\in C_t}\sgn(c)c\{t\}
\]
maps to $v_L$: in a column, the position action replaces a permutation
of the letters by its inverse, which has the same sign.
In particular this vector is nonzero, since the original row word has
coefficient $1$.  The span of all translates of $e_t$ is the Specht
module $S^\theta$, irreducible over $\C$ by the standard polytabloid
construction \cite[Definitions~4.1 and~4.3, Theorem~4.12]{james1978}.
This proves the assertion.
The argument applies to a dual space with any chosen basis; alternatively,
an invertible basis change on that space induces an invertible map on
its position tensor power commuting with $S_B$.
\end{proof}

We use the usual complex representation-theoretic facts
\cite[Theorems~4.12 and~6.7]{james1978}: the $S^\lambda$, for
$\lambda\vdash n$, form the complete list of irreducible $S_n$-modules,
and
\begin{equation}\label{tools:dual-sign}
 (S^\lambda)^*\cong S^\lambda,
 \qquad S^\lambda\otimes\sgn\cong S^{\lambda^t}.
\end{equation}
Every complex representation of a finite group is semisimple
\cite[Proposition~1.5 and Corollary~1.6]{fultonharris1991}, so any
irreducible quotient is also a constituent.  For $H\le G$,
Frobenius reciprocity \cite[Proposition~3.17]{fultonharris1991}
takes the form
\begin{equation}\label{tools:frobenius}
 \Hom_G(\Ind_H^G A,V)
 \cong\Hom_H(A,\Res_H^G V).
\end{equation}
Induction, modeled by $\C[G]\otimes_{\C[H]}A$, is transitive and
distributes over direct sums.  The symbol
$\boxtimes$ denotes the outer tensor product for separate group actions;
ordinary tensor products of two $S_B$-modules carry the diagonal action.
The natural pairing of $E^{\otimes B}$ and $(E^*)^{\otimes B}$ satisfies
$\langle gu,gt\rangle=\langle u,t\rangle$.
Thus pairing with an invariant subspace $T$ of the dual tensor space
defines an equivariant map $u\mapsto(t\mapsto\langle u,t\rangle)$
into $T^*$.

\subsection{Pieri and the needed part of Young's rule}

For diagrams $\nu\subseteq\lambda$, the skew diagram $\lambda/\nu$
is a \emph{horizontal strip} if it has at most one box in each column.
We use the empty partition and the trivial group $S_0$ in the following
standard identity.

\begin{lemma}[Pieri rule]\label{lem:pieri}
For $\nu\vdash a$ and $b\ge0$,
\begin{equation}\label{tools:pieri}
 \Ind_{S_a\times S_b}^{S_{a+b}}
      (S^\nu\boxtimes\mathbf1)
 \cong
 \bigoplus_{\substack{\lambda\vdash a+b,\ \lambda\supseteq\nu\\
                      \lambda/\nu\text{ a horizontal strip}}}
       S^\lambda.
\end{equation}
Every summand occurs once.
\end{lemma}

This is the one-row specialization of the Littlewood--Richardson rule
\cite[Section~16, especially Rule~16.4]{james1978}: all $b$ entries of
its skew filling are $1$, so column strictness requires a horizontal
strip; the unique such filling satisfies the lattice-word condition.
We give the dominance and multiplicity
consequences needed here, including the constructive existence argument.

\begin{lemma}[Young's rule: support and diagonal multiplicity]
\label{lem:young-rule}
Let $\theta=(\theta_1,\ldots,\theta_\ell)\vdash n$. The row-tabloid
permutation module has the equivalent induced realization
\[
 M^\theta=
 \Ind_{S_{\theta_1}\times\cdots\times S_{\theta_\ell}}^{S_n}
 \mathbf1.
\]
Then $S^\tau$ occurs in $M^\theta$ if and only if
$\tau\unrhd\theta$, and $S^\theta$ occurs once.
The word module of any content obtained by reordering the parts of
$\theta$ is also isomorphic to $M^\theta$.
\end{lemma}

\begin{proof}
The case $n=0$ is immediate from the trivial representation of $S_0$.
Assume henceforth that $n>0$.
Transitivity of induction and \Cref{lem:pieri} show that the multiplicity
of $S^\tau$ counts chains
\begin{equation}\label{tools:strip-chain}
 \varnothing=\nu^{(0)}\subseteq\nu^{(1)}\subseteq\cdots
 \subseteq\nu^{(\ell)}=\tau,
 \qquad |\nu^{(k)}/\nu^{(k-1)}|=\theta_k,
\end{equation}
whose successive differences are horizontal strips.
Each strip raises the first column's height by at most one, so
$\nu^{(k)}$ has at most $k$ rows.  Since it is contained in $\tau$,
\[
 \sum_{i=1}^k\theta_i=|\nu^{(k)}|
 \le\sum_{i=1}^k\tau_i.
\]
For $k>\ell$, both initial sums equal $n$, since $\tau$ has at most
$\ell$ rows.  This proves necessity of dominance.
If $\tau=\theta$, equality forces
$\nu^{(k)}$ to be exactly the first $k$ rows of $\theta$.
These diagrams do form a chain of the required kind, proving
multiplicity one.

For sufficiency, suppose $\tau\unrhd\theta$ and induct on $\ell$.
For $\ell=1$, necessarily $\tau=(n)$.
For $\ell>1$, put $t=\theta_\ell$.
There are at least $t$ columns of $\tau$, since
$\tau_1\ge\theta_1\ge t$.  Delete a bottom box from each of the
$t$ tallest columns, choosing the rightmost columns first among equal
heights, and call the remaining diagram $\nu$.
Its column heights are still nonincreasing: unequal adjacent heights
differ by at least one, and within any group of equal heights precisely
a suffix is lowered.  Thus $\nu$ is a partition and $\tau/\nu$ is a
horizontal $t$-strip.  For every $k\ge1$ its loss in the first $k$ rows is
\begin{equation}\label{tools:prefix-loss}
 \sum_{i=1}^k\tau_i-\sum_{i=1}^k\nu_i
       =\max(0,t-\tau_{k+1}):
\end{equation}
there are $\tau_{k+1}$ columns taller than $k$, and the rule selects these
before any column whose bottom box belongs to the first $k$ rows.
If the loss is zero, the desired dominance inequality for $\nu$ follows
immediately.  If it is positive and $k<\ell$, then
\[
 \sum_{i=1}^k\nu_i
 =\sum_{i=1}^{k+1}\tau_i-t
 \ge\sum_{i=1}^{k+1}\theta_i-t
 \ge\sum_{i=1}^k\theta_i.
\]
For completeness, dominance at $\ell$ gives $\len(\tau)\le\ell$,
and dominance at $\ell-1$ gives $\tau_\ell\le t$.
All $\tau_\ell$ columns of height $\ell$ are therefore selected, so
the last row disappears and $\len(\nu)\le\ell-1$.
Hence all remaining initial sums equal the total $n-t$, and
$\nu\unrhd(\theta_1,\ldots,\theta_{\ell-1})$.
Induction gives a chain ending at $\nu$; appending the deleted strip
gives \Cref{tools:strip-chain} ending at $\tau$.

Finally, the stabilizer of a word with prescribed letter multiplicities
is the product of the symmetric groups on its equal-letter position
sets.  Reordering those multiplicities conjugates this subgroup, giving
the asserted isomorphism of permutation modules.
\end{proof}

\section{The triangular cyclic module and dominance}
\label{sec:dominance}

We now fix the generator from the introduction exactly and state the
cyclic support theorem used in the induction. The second task of this
section is to detect every constituent dominated by the staircase,
using a sign twist and an explicit permutation-module quotient.

Following the triangular arrangement in \cite[Section~4]{ikenmeyer2015},
for $m\ge1$ use the position set
\begin{equation}
 B_m=\{(i,j)\in\mathbb Z_{\ge1}^2:i+j\le m+1\},
 \qquad G_m=S_{B_m}.
 \label{dom:triangle}
\end{equation}
Its row and column set partitions are
\[
 R_m=\{R_{m,i}:1\le i\le m\},\qquad
 R_{m,i}=\{(i,j):1\le j\le m+1-i\},
\]
\[
 C_m=\{C_{m,j}:1\le j\le m\},\qquad
 C_{m,j}=\{(i,j):1\le i\le m+1-j\}.
\]
Order each row by increasing $j$ and each column by increasing $i$.
With $E_m=\C^m$ and its ordered basis $e_1,\ldots,e_m$, form the
alternating block tensors of \Cref{lem:polytabloid} and set
\begin{equation}
 w_m=v_{R_m}\otimes v_{C_m},\qquad
 W_m=\C[G_m]w_m
 \ \subseteq\ E_m^{\otimes B_m}\otimes E_m^{\otimes B_m}.
 \label{dom:cyclic-module}
\end{equation}
Both set partitions have block sizes $m,m-1,\ldots,1$.
Since $\rho_m^t=\rho_m$, \Cref{lem:polytabloid} gives
\begin{equation}
 V_R:=\C[G_m]v_{R_m}\cong S^{\rho_m},\qquad
 V_C:=\C[G_m]v_{C_m}\cong S^{\rho_m},\qquad
 W_m\subseteq V_R\otimes V_C.
 \label{dom:square-inclusion}
\end{equation}
Here and throughout, the action on a tensor product of two $G_m$-modules
is diagonal.  Thus it suffices to prove the following stronger statement.

\begin{theorem}[Cyclic staircase support]
\label{thm:cyclic-saxl}
For every $m\ge1$ and every partition $\lambda\vdash N_m$, the irreducible
$S^\lambda$ occurs in $W_m$.
\end{theorem}

We prove \Cref{thm:cyclic-saxl} by induction in \Cref{sec:completion}.
The present section establishes the case supplied by dominance.

\begin{proposition}
\label{prop:dominated}
If $m\ge1$, $\lambda\vdash N_m$, and $\rho_m\unrhd\lambda$, then
$S^\lambda$ occurs in $W_m$.
\end{proposition}

\begin{proof}
Fix $m$, abbreviate $G=G_m$, and let $\epsilon$ be the one-dimensional
sign representation, also denoting its character by $\epsilon$.
By Young's rule, conjugation of dominance, and the sign-twist identity
in \Cref{tools:dual-sign}, it is enough to show that
$W_m\otimes\epsilon$ contains every constituent of $M^{\rho_m}$.
We will find the cyclic module generated by
$v_{R_m}\otimes v_{R_m}$ inside an isomorphic image of
$W_m\otimes\epsilon$, then map that module onto $M^{\rho_m}$.
The row subgroup
\[
 H_R=\prod_{i=1}^m S_{R_{m,i}}
\]
acts on $v_{R_m}$ by $\epsilon|_{H_R}$.  Consequently the unnormalized signed
diagonal average is exactly
\begin{equation}
 \sum_{h\in H_R}\epsilon(h)\,h w_m
 =v_{R_m}\otimes u,\qquad
 u=\sum_{h\in H_R}h v_{C_m}.
 \label{dom:signed-average}
\end{equation}
To see that $u\ne0$, let $a_0(i,j)=i$ and consider the corresponding word
$e_{a_0}$.  It is fixed by $H_R$.  Its coefficient in $v_{C_m}$ is $1$:
within each increasingly ordered column it is the identity assignment
$1,2,\ldots,m+1-j$.  The coefficient of $e_{a_0}$ in every $h v_{C_m}$ is
therefore also $1$, since $h^{-1}e_{a_0}=e_{a_0}$.  Its coefficient in $u$ is
$|H_R|=\prod_{r=1}^m r!$, which is nonzero.

We next identify the alternating line in $V_R$.  Frobenius reciprocity,
the induction identity for a sign twist, and \Cref{lem:young-rule} give
\begin{align}
 \dim\Hom_{H_R}(\epsilon|_{H_R},\Res_{H_R}^{G}V_R)
 &=\dim\Hom_G(\Ind_{H_R}^{G}\epsilon|_{H_R},V_R)\notag\\
 &=\dim\Hom_G(M^{\rho_m}\otimes\epsilon,V_R)\notag\\
 &=\dim\Hom_G(M^{\rho_m},V_R\otimes\epsilon)=1.
 \label{dom:alternating-line}
\end{align}
For the last equality, $V_R\otimes\epsilon\cong S^{\rho_m}$ by
self-conjugacy of the staircase, and $S^{\rho_m}$ has multiplicity one in
$M^{\rho_m}$.  As $v_{R_m}$ is nonzero and alternating under $H_R$, it spans
the line in \Cref{dom:alternating-line}.

Twisting the diagonal tensor product can be done in its second factor:
the reassociation
\begin{equation}
 (V_R\otimes V_C)\otimes\epsilon
 \longrightarrow V_R\otimes(V_C\otimes\epsilon),\qquad
 (x\otimes y)\otimes z\longmapsto x\otimes(y\otimes z)
 \label{dom:twist-one-factor}
\end{equation}
is a $G$-isomorphism.  Choose a $G$-isomorphism
$\phi:V_C\otimes\epsilon\longrightarrow V_R$, which exists by
\Cref{dom:square-inclusion} and $\rho_m^t=\rho_m$, and choose $0\ne z\in\epsilon$.
Since $u$ is $H_R$-fixed, $u\otimes z$ is alternating under $H_R$.
It follows from \Cref{dom:alternating-line} that
$\phi(u\otimes z)=c\,v_{R_m}$ for some $c\ne0$.
Tensor \Cref{dom:signed-average} with $z$ and apply the reassociation in
\Cref{dom:twist-one-factor}, followed by $\id\otimes\phi$.
The right-hand side becomes $c(v_{R_m}\otimes v_{R_m})$.
Thus the image of $W_m\otimes\epsilon$ under this isomorphism contains
\begin{equation}
 Z_m=\C[G](v_{R_m}\otimes v_{R_m})
 \ \subseteq\ V_R\otimes V_R.
 \label{dom:repeated-row-cyclic}
\end{equation}

We construct a surjection from $Z_m$ onto $M^{\rho_m}$.
Let $E_{\mathrm{out}}$ have basis $f_1,\ldots,f_{2m-1}$ and define the linear map
\[
 q:E_m\otimes E_m\longrightarrow E_{\mathrm{out}},\qquad
 q(e_a\otimes e_b)=f_{a+b-1}.
\]
Regroup the two input layers by position, using ordinary tensor
reassociation, and apply $q$ at every position.  Then project onto the
span $\mathcal M_\alpha$ of words in which letter $r$ occurs exactly $r$
times for $1\le r\le m$, and no letter greater than $m$ occurs.
Thus the content is $\alpha=(1,2,\ldots,m,0,\ldots,0)$ in the output
alphabet.  Denote the resulting linear map by
\begin{equation}
 Q:E_m^{\otimes B_m}\otimes E_m^{\otimes B_m}
 \longrightarrow\mathcal M_\alpha.
 \label{dom:pair-letter-map}
\end{equation}
Both operations commute with position permutations: the same local map
$q$ is used everywhere, and content is unchanged by $G$.  Hence $Q$ is
$G$-equivariant.

Consider any summand in the expansion of $v_{R_m}\otimes v_{R_m}$.
On a row of length $r$, write its two assignments as
$(a_1,\ldots,a_r)$ and $(b_1,\ldots,b_r)$, each a permutation of
$1,\ldots,r$, and put $c_j=a_j+b_j-1$.  Then
\begin{equation}
 \sum_{j=1}^r c_j=r^2,
 \qquad
 \sum_{j=1}^r c_j^2-r^3
 =\sum_{j=1}^r(c_j-r)^2\ge0.
 \label{dom:square-minimum}
\end{equation}
For a term retained by the content projection, the sum of squared output
letters over the whole triangle is $\sum_{r=1}^m r^3$.
Summing \Cref{dom:square-minimum} over the rows forces equality on every
row.  Thus $c_j=r$ at every position in the row of length $r$.
There is only one possible retained output word, namely
\[
 t_m=\bigotimes_{(i,j)\in B_m}f_{m+1-i}.
\]
It remains to compute its coefficient.  On a row of length $r$, each of
the $r!$ assignments $a_j=\pi(j)$ has the unique compatible assignment
$b_j=r+1-\pi(j)$.  The latter is the permutation $\omega_r\pi$, where
$\omega_r(j)=r+1-j$, so the product of the two alternating signs is
\[
 \sgn(\pi)\sgn(\omega_r\pi)
 =\sgn(\omega_r)=(-1)^{\binom r2}.
\]
The row contributions are independent.  In particular, there is no
cancellation, and the exact image is
\begin{equation}
 Q(v_{R_m}\otimes v_{R_m})
 =\left((-1)^{\sum_{r=1}^m\binom r2}\prod_{r=1}^m r!\right)t_m\ne0.
 \label{dom:noncancellation}
\end{equation}
The stabilizer of $t_m$ is exactly $H_R$, since its equal-letter sets are
the rows of $B_m$.  Its orbit is a basis of $\mathcal M_\alpha$, and hence
\begin{equation}
 Q(Z_m)=\C[G]t_m=\mathcal M_\alpha
 \cong\Ind_{H_R}^{G}\mathbf1\cong M^{\rho_m}.
 \label{dom:permutation-quotient}
\end{equation}
This also explains why the increasing content $(1,2,\ldots,m)$ gives the
usual module indexed by the decreasing partition $\rho_m$: reordering
the content parts conjugates the standard Young subgroup.

By \Cref{lem:young-rule}, every $S^\tau$ with $\tau\unrhd\rho_m$ occurs in
the quotient in \Cref{dom:permutation-quotient}.  Semisimplicity and
\Cref{dom:repeated-row-cyclic} imply that it occurs in $W_m\otimes\epsilon$.
If $\rho_m\unrhd\lambda$, reversal of dominance under conjugation gives
$\lambda^t\unrhd\rho_m^t=\rho_m$.  Taking $\tau=\lambda^t$ and twisting
back by $\epsilon$ proves that $S^\lambda$ occurs in $W_m$.
\end{proof}

\section{A quotient obtained by cutting off a band}
\label{sec:band-quotient}

The next step transfers support from a smaller cyclic staircase module
and a band into $W_m$. The projection must retain their specified
generators and realize the full induced module, not just a submodule
of it. We first define the two factors and then verify both requirements.

Fix $s\in\{1,2\}$ and $h=m-s\ge1$, and put
\begin{equation}\label{band:sets}
 D=B_h,\qquad F=B_m\setminus D,\qquad
 G=G_m=S_{B_m},\qquad H=S_D\times S_F.
\end{equation}
Thus $|F|=N_m-N_h=s(2m-s+1)/2$.
Let $R_F$ and $C_F$ be the partitions of $F$ into its nonempty row
and column intersections, with the orders inherited from $R_m$ and $C_m$.
Each block has size at most $s$.  In two layers with alphabet
$1,\ldots,s$, define
\begin{equation}\label{band:generator}
 u_{m,s}=v_{R_F}\otimes v_{C_F},\qquad
 U_{m,s}=\C[S_F]u_{m,s}.
\end{equation}
The dependence on $m$ records the length of the band.
Our goal is a quotient of $W_m$ induced from $W_h\boxtimes U_{m,s}$.

In the disjoint-coset model of induction
\cite[Section~3.3]{fultonharris1991}, translates of the inducing module
form a direct sum. We record why disjoint coordinate sectors realize
this model for the cyclic span of a specified vector, including the
inverse map.

\begin{lemma}[Induction from coordinate sectors]\label{lem:sector-induction}
Let a finite group $G$ act transitively on a finite set $\Omega$.
Suppose a $G$-module has a direct-sum decomposition
$T=\bigoplus_{A\in\Omega}T_A$ with $gT_A=T_{gA}$.
Fix $D\in\Omega$, let $H$ be its stabilizer, and take $z\in T_D$.
If $A_0=\C[H]z$, then the map
\begin{equation}\label{band:induced-map}
 \Phi:\C[G]\otimes_{\C[H]}A_0\longrightarrow\C[G]z,
 \qquad g\otimes a\longmapsto ga
\end{equation}
is an isomorphism of $G$-modules.
\end{lemma}

\begin{proof}
The relation $(gh)\otimes a=g\otimes ha$, for $h\in H$, has the
same image on both sides, so $\Phi$ is well-defined.  It is equivariant
and surjective.  Choose $g_A\in G$ with $g_AD=A$ for every $A\in\Omega$,
and let $\pi_A:T\to T_A$ be the coordinate projection.  Since
\[
 \C[G]z=\bigoplus_{A\in\Omega}g_AA_0,
 \qquad g_AA_0\subseteq T_A,
\]
we have $g_A^{-1}\pi_A(y)\in A_0$ for $y\in\C[G]z$.  Define
\begin{equation}\label{band:induced-inverse}
 \Psi(y)=\sum_{A\in\Omega}
          g_A\otimes g_A^{-1}\pi_A(y).
\end{equation}
Replacing $g_A$ by $g_Ah$ does not change this expression, because
$g_Ah\otimes h^{-1}a=g_A\otimes a$ in the balanced tensor product.
Clearly $\Phi\Psi(y)=\sum_A\pi_A(y)=y$.
For $g\otimes a$, set $A=gD$ and write $g=g_Ah$ with $h\in H$.
Only its $A$-sector is nonzero, and hence
\[
 \Psi\Phi(g\otimes a)=g_A\otimes ha=g\otimes a.
\]
Thus the displayed maps are inverse isomorphisms.
\end{proof}

\begin{proposition}[Band quotient]\label{prop:band-quotient}
For the parameters in \Cref{band:sets}, there is a surjective
$G_m$-equivariant map
\begin{equation}\label{band:quotient}
 W_m\twoheadrightarrow
 \Ind_{S_D\times S_F}^{S_{B_m}}(W_h\boxtimes U_{m,s}).
\end{equation}
In particular, every irreducible constituent of the induced module on
the right is a constituent of $W_m$.
\end{proposition}

\begin{proof}
\emph{Projection and the high positions.}
Use the ordinary tensor identification
\[
 E_m^{\otimes B_m}\otimes E_m^{\otimes B_m}
 \cong (E_m\otimes E_m)^{\otimes B_m}.
\]
Put $E_{\mathrm{lo}}=\Span(e_1,\ldots,e_s)$ and
$E_{\mathrm{hi}}=\Span(e_{s+1},\ldots,e_m)$.
On a pair-letter basis vector define
\[
 p_s(e_a\otimes e_b)=
 \begin{cases}
 e_a\otimes e_b,& a,b\le s\text{ or }a,b>s,\\
 0,&\text{otherwise},
 \end{cases}
 \qquad P_s=p_s^{\otimes B_m}.
\]
The same local projection is used at every position, so $P_s$
commutes with $G$.  Call a retained position \emph{high} when both
letters are greater than $s$.

In a retained word of $w_m$, let $A$ be the set of high positions.
Alternation in the first layer gives row counts
$r_i=\max(h+1-i,0)$, and alternation in the second gives column
counts $d_j=\max(h+1-j,0)$, for $1\le i,j\le m$.
For each $1\le k\le m$, these counts satisfy
\begin{equation}\label{band:margin-equality}
 \sum_{i=1}^k r_i
 =\#\{(i,j)\in B_h:i\le k\}
 =\sum_{j=1}^m\min(k,d_j).
\end{equation}
Column $j$ can contribute at most $\min(k,d_j)$ high positions to
the first $k$ rows.  Equality of the sums in
\Cref{band:margin-equality} therefore forces equality in every column.
For $d_j>0$, take $k=d_j$: all its $d_j$ high positions occupy the
first $d_j$ rows.  Columns with $d_j=0$ have no high positions.
Consequently
\begin{equation}\label{band:unique-high-set}
 A=\{(i,j)\in B_m:i\le h+1-j\}=B_h=D.
\end{equation}
For the zero-one matrix recording membership in $A$, this is the
equality case of the row and column margin inequalities
\cite[Section~2]{ryser1957}. The specific staircase margins in
\Cref{band:margin-equality} make every prefix bound an equality.

\emph{Factorization with signs.}
In a row or column block of length $r$, its intersection with $D$
is its initial segment of length $a=\max(r-s,0)$; its intersection
with $F$ has length $b=\min(r,s)$.
The retained assignments put the high letters $s+1,\ldots,r$ on
the first segment and the low letters $1,\ldots,b$ on the second.
Let $\Alt$ denote the unnormalized alternation of a displayed list,
placed in the indicated segment, with $\Alt(\varnothing)=1$.
The part of the block tensor with this prescribed high set is exactly
\begin{equation}\label{band:block-factorization}
 (-1)^{ab}\,
 \Alt(e_{s+1},\ldots,e_r)\otimes\Alt(e_1,\ldots,e_b).
\end{equation}
Here the high list is empty when $r\le s$.
Indeed, every high letter precedes every low letter in the position
order, giving $ab$ inversions; all other inversions are internal to
the two independently permuted lists.  Every such pair of permutations
occurs once.  In particular, when $r<s$ the formula is just the
original alternation on $e_1,\ldots,e_r$.

Rows and columns both have the length list $m,m-1,\ldots,1$.
Thus their cross signs in \Cref{band:block-factorization} multiply to
\[
 \prod_{r=1}^m(-1)^{2\max(r-s,0)\min(r,s)}=1.
\]
Identify $E_h$ with $E_{\mathrm{hi}}$ by
$e_a\mapsto e_{s+a}$, and write $\widetilde w_h$ and
$\widetilde W_h$ for the images of $w_h$ and $W_h$ in the two
high-letter layers on $D$.  The low portions are precisely the blocks
of $R_F,C_F$.  Regrouping the factors over the disjoint sets $D,F$
introduces no signs, and hence
\begin{equation}\label{band:generator-factorization}
 z:=P_s(w_m)=\widetilde w_h\otimes u_{m,s}\ne0.
\end{equation}
The nonvanishing follows also directly from the nonzero alternating
tensors on each disjoint block.

The two factors of $H=S_D\times S_F$ act independently on this
tensor, while each acts diagonally on the two layers over its own
position set.  Therefore
\begin{align}
 A_0:=\C[H]z
 &=\Span\{(\sigma\widetilde w_h)\otimes(\tau u_{m,s}):
             \sigma\in S_D,\ \tau\in S_F\}\notag\\
 &=\widetilde W_h\boxtimes U_{m,s}
 \cong W_h\boxtimes U_{m,s}.
 \label{band:subgroup-orbit}
\end{align}
The second equality follows by expanding tensors of linear
combinations of the two separate orbit sets.

\emph{The coordinate sectors.}
Let $\Omega$ be the set of subsets $A\subset B_m$ of size $|D|$.
Within the pair-letter space define
\[
 T_A=(E_{\mathrm{hi}}\otimes E_{\mathrm{hi}})^{\otimes A}
       \otimes
       (E_{\mathrm{lo}}\otimes E_{\mathrm{lo}})^{\otimes(B_m\setminus A)}.
\]
These are the coordinate spans of words high exactly on $A$.
They have disjoint word bases, so their sum is direct, and
$gT_A=T_{gA}$.  By \Cref{band:generator-factorization}, $z\in T_D$;
the stabilizer of $D$ in $G$ is exactly $H$.
Apply \Cref{lem:sector-induction} and
\Cref{band:subgroup-orbit} to obtain
\[
 \C[G]z\cong\Ind_H^G A_0
       \cong\Ind_H^G(W_h\boxtimes U_{m,s}).
\]
Since $P_s$ is equivariant, $P_s(W_m)=\C[G]P_s(w_m)=\C[G]z$.
Its restriction to $W_m$, followed by the inverse map $\Psi$ in
\Cref{band:induced-inverse} and the identification of $A_0$ in
\Cref{band:subgroup-orbit}, gives \Cref{band:quotient}.
The image of this ambient projection need not lie in the two original
Specht factors; the composition is still a surjective $G$-equivariant
linear map onto
the displayed module.  Semisimplicity over $\C$ proves the last
assertion.
\end{proof}

\section{Cyclic support of the bands}\label{sec:path-band}

For the band step of the induction, it remains to determine the support
needed in $U_{m,1}$ and $U_{m,2}$ for \Cref{prop:band-quotient}.
The width-one factor is trivial.
For width two we will detect every shape with at most four rows directly
in the specified generator, using the invariant evaluation pairing.

For width one, the band consists of the $m$ positions with $i+j=m+1$.
Every row and column block is a singleton, and both alphabets consist only
of the letter $1$. Thus $u_{m,1}$ is a nonzero constant word tensor, fixed
by every position permutation, and
\begin{equation}\label{path:one-band}
 U_{m,1}\cong S^{(m)} \qquad (m\geq2).
\end{equation}
We now determine the support needed from the particular width-two
generator.

\subsection{The indexed path contraction}

Let $m\geq3$, $K=2m-1$, and $F=B_m\setminus B_{m-2}$.
Number its positions along the path by
\begin{equation}\label{path:positions}
 p_{2j-1}=(m+1-j,j)\quad(1\leq j\leq m),\qquad
 p_{2j}=(m-j,j)\quad(1\leq j<m).
\end{equation}
In this numbering, the row blocks are the singleton $1$ and the pairs
$(2,3),(4,5),\ldots,(K-1,K)$; the column blocks are
$(1,2),(3,4),\ldots,(K-2,K-1)$ and the singleton $K$.
Order each pair by increasing path index and let $u$ denote the resulting
alternating band tensor. The row orders agree with increasing $j$, while
each of the $m-1$ column pairs reverses increasing $i$. Consequently
\begin{equation}\label{path:generator-sign}
 u=(-1)^{m-1}u_{m,2},\qquad U_{m,2}=\C[S_F]u.
\end{equation}
\Cref{fig:band-path} illustrates the numbering and the consecutive
segments used in \Cref{path:conjugation-example}.

\begin{figure}[ht]
\centering
\setlength{\unitlength}{1pt}
\begin{picture}(430,155)
\small
\put(90,148){\makebox(0,0){Band in $B_3$}}
\put(10,139){\vector(1,0){25}}
\put(39,139){\makebox(0,0){$j$}}
\put(10,139){\vector(0,-1){25}}
\put(10,108){\makebox(0,0){$i$}}
\put(45,110){\color[gray]{0.65}\circle*{6}}
\put(45,123){\makebox(0,0){$D=B_1$}}
\put(45,30){\circle*{4}}
\put(45,70){\circle*{4}}
\put(90,70){\circle*{4}}
\put(90,110){\circle*{4}}
\put(135,110){\circle*{4}}
\put(45,34){\vector(0,1){32}}
\put(90,74){\vector(0,1){32}}
\multiput(49,70)(8,0){4}{\line(1,0){4}}
\put(81,70){\vector(1,0){5}}
\multiput(94,110)(8,0){4}{\line(1,0){4}}
\put(126,110){\vector(1,0){5}}
\put(45,16){\makebox(0,0){$p_1$}}
\put(30,70){\makebox(0,0){$p_2$}}
\put(106,70){\makebox(0,0){$p_3$}}
\put(90,124){\makebox(0,0){$p_4$}}
\put(135,124){\makebox(0,0){$p_5$}}
\put(305,148){\makebox(0,0){Dual-test columns for $\eta=(2,2,1)$}}
\multiput(220,90)(40,0){3}{\circle*{4}}
\put(350,90){\circle*{4}}
\put(390,90){\circle*{4}}
\put(224,90){\vector(1,0){32}}
\multiput(264,90)(8,0){3}{\line(1,0){4}}
\put(288,90){\vector(1,0){8}}
\put(304,90){\vector(1,0){42}}
\multiput(354,90)(8,0){3}{\line(1,0){4}}
\put(378,90){\vector(1,0){8}}
\put(220,76){\makebox(0,0){$p_1$}}
\put(260,76){\makebox(0,0){$p_2$}}
\put(300,76){\makebox(0,0){$p_3$}}
\put(350,76){\makebox(0,0){$p_4$}}
\put(390,76){\makebox(0,0){$p_5$}}
\put(325,70){\color[gray]{0.5}\line(0,1){38}}
\put(215,55){\line(1,0){90}}
\put(215,55){\line(0,1){6}}
\put(305,55){\line(0,1){6}}
\put(260,43){\makebox(0,0){$r_1=3$}}
\put(335,55){\line(1,0){60}}
\put(335,55){\line(0,1){6}}
\put(395,55){\line(0,1){6}}
\put(365,43){\makebox(0,0){$r_2=2$}}
\end{picture}
\caption{The width-two band for $m=3$. Solid arrows are column pairs;
dashed arrows are row pairs, all directed by increasing path index.
Column arrows reverse increasing $i$, as in
\Cref{path:generator-sign}. The row singleton is $p_1$ and the column
singleton is $p_5$, each carrying the fixed letter $1$ in its layer.
On the right, the dual-test tableau of shape $\eta=(2,2,1)$ has column
lengths $r_1=3$ and $r_2=2$, giving the two bracketed consecutive segments.
The gray vertical line separates these segments, which differ from the
original pair blocks.}
\label{fig:band-path}
\end{figure}

Put $E=\C^2$ with basis $e_1,e_2$ and $V=E\otimes E$.
Ordinary tensor reassociation identifies the two band layers with
$V^{\otimes K}$, using the position order in \Cref{path:positions}.
Identify $V^*$ with $\Mat_2(\C)$ by letting a matrix $X$ take the value
$X_{ab}$ on $e_a\otimes e_b$. All pairings below are bilinear evaluation
pairings. Write
\[
 J=\begin{pmatrix}0&1\\-1&0\end{pmatrix}.
\]

\begin{lemma}[Path contraction]\label{path:contraction}
For arbitrary $X_1,\ldots,X_K\in\Mat_2(\C)$,
\begin{equation}\label{path:contraction-formula}
 \left\langle u,X_1\otimes\cdots\otimes X_K\right\rangle
 =(-1)^{m-1}
 \left[X_1\adj(X_2)X_3\adj(X_4)\cdots
                 X_{K-2}\adj(X_{K-1})X_K\right]_{11}.
\end{equation}
\end{lemma}

\begin{proof}
The two singleton blocks require the first row-layer index and last
column-layer index to be $1$. The coefficient of an alternating pair in
path order is $J_{ab}$. Hence the left side of
\Cref{path:contraction-formula} is exactly
\begin{equation}\label{path:index-sum}
 \sum_{\substack{a_1,\ldots,a_K,b_1,\ldots,b_K\in\{1,2\}\\
                  a_1=b_K=1}}
 \left(\prod_{\ell=1}^{m-1}
   J_{b_{2\ell-1},b_{2\ell}}J_{a_{2\ell},a_{2\ell+1}}\right)
 \prod_{q=1}^K(X_q)_{a_qb_q}.
\end{equation}
Reading the contracted indices successively from $a_1$ to $b_K$ gives
\[
 \left[X_1JX_2^{\mathsf T}JX_3JX_4^{\mathsf T}J
       \cdots X_{K-2}JX_{K-1}^{\mathsf T}JX_K\right]_{11}.
\]
The transpose at an even position records that its column index is
encountered before its row index. For
$X=\left(\begin{smallmatrix}a&b\\c&d\end{smallmatrix}\right)$,
direct multiplication gives
\begin{equation}\label{path:adjugate}
 \adj(X)=\begin{pmatrix}d&-b\\-c&a\end{pmatrix}
        =\tr(X)I-X,\qquad JX^{\mathsf T}J=-\adj(X).
\end{equation}
There are $m-1$ even positions, yielding the stated sign.
In particular, adjugation is a linear operation on $\Mat_2(\C)$.
\end{proof}

\subsection{Alternating consecutive segments}

For a shape with at most four rows, place its tableau columns on
successive disjoint segments of the path. On a segment of length $r$,
the dual polytabloid alternates the first $r$ elements of one common
basis of $V^*$. The contraction formula therefore assigns an alternating
matrix sum to each segment; these sums multiply in path order.
We now choose the basis so that every such segment matrix, for
$1\le r\le4$ and either starting parity, is invertible.

Use the following ordered basis of $V^*$:
\begin{equation}\label{path:matrix-basis}
 M_1=I,\qquad
 M_2=Z=\begin{pmatrix}1&0\\0&-1\end{pmatrix},\qquad
 M_3=T=\begin{pmatrix}0&1\\1&0\end{pmatrix},\qquad
 M_4=J=ZT.
\end{equation}
These matrices are linearly independent: $I,Z$ span the diagonal
matrices and $T,J$ span the off-diagonal matrices. The last three are
traceless and anticommute in pairs; moreover
\[
 Z^2=T^2=I,\qquad J^2=-I.
\]
Thus all four are invertible. Alternating evaluations on $Z,T,J$ are also
used to test polynomial identities for two-by-two matrices with adjugation
\cite[Section~3]{drenskygiambruno1994}; here the adjugations are determined
by the path positions.

For a global path position $q$, define the linear operation
\[
 D_q(X)=
 \begin{cases}
 X,&q\text{ odd},\\
 \adj(X),&q\text{ even}.
 \end{cases}
\]
For a segment of length $r\leq4$ beginning at position $p$, define its
alternating matrix by the following sum of ordered products:
\begin{equation}\label{path:segment-definition}
 A_{r,p}=\sum_{\pi\in S_r}\sgn(\pi)
      D_p(M_{\pi(1)})D_{p+1}(M_{\pi(2)})\cdots
      D_{p+r-1}(M_{\pi(r)}).
\end{equation}

\begin{lemma}[Segment alternation]\label{path:segment}
Let $e$ be the number of even positions among $p,\ldots,p+r-1$, and put
$\delta=1$ if $p$ is even and $\delta=0$ otherwise. Then
\begin{equation}\label{path:segment-formula}
 A_{r,p}=r!(-1)^{e-\delta}M_1\cdots M_r.
\end{equation}
In particular, all these matrices are invertible. Their values are
\begin{equation}\label{path:segment-table}
 \begin{array}{c|cc}
 r & p\text{ odd} & p\text{ even}\\ \hline
 1&I&I\\
 2&-2Z&2Z\\
 3&-6J&-6J\\
 4&-24I&24I
 \end{array}
\end{equation}
\end{lemma}

\begin{proof}
Fix a permutation $\pi$ and let $j$ be the local slot occupied by $M_1=I$.
Deleting $I$ leaves a permutation of $M_2,\ldots,M_r$.
Its sign is exactly the sign acquired when those pairwise anticommuting
matrices are reordered. Moving $I$ from its first canonical slot to slot
$j$ contributes $(-1)^{j-1}$ to $\sgn(\pi)$ and changes no matrix product.
It follows that
\[
 \sgn(\pi)M_{\pi(1)}\cdots M_{\pi(r)}
       =(-1)^{j-1}M_1\cdots M_r.
\]
By \Cref{path:adjugate}, adjugation fixes $I$ and negates each of the
other basis matrices. Write $\epsilon_j=1$ if $p+j-1$ is even, and $0$
otherwise. Adjugation therefore contributes the additional sign
$(-1)^{e-\epsilon_j}$. Since
$\epsilon_j\equiv\delta+j-1\pmod2$, the total sign is
$(-1)^{e-\delta}$, independently of $\pi$. All $r!$ summands in
\Cref{path:segment-definition} consequently agree, proving
\Cref{path:segment-formula}, including $r=1$.
Finally, the canonical products for $r=1,2,3,4$ are respectively
$I,Z,J,-I$, which gives \Cref{path:segment-table}.
\end{proof}

\subsection{Detection by a dual polytabloid}

The segment calculation gives an invertible ordered product. Its fixed
endpoint entry can nevertheless vanish. We now choose the dual test
basis so that this entry is nonzero, while keeping the band generator
and its endpoint letters unchanged.

\begin{proposition}[Four-row support of the width-two band]\label{prop:two-band}
For every $m\geq3$ and every partition $\eta\vdash 2m-1$ with
$\len(\eta)\leq4$, the irreducible module $S^\eta$ occurs in $U_{m,2}$.
\end{proposition}

\begin{proof}
Let $r_1,\ldots,r_c$ be the column lengths of $\eta$ in their usual order,
so $1\leq r_j\leq4$ and $\sum_jr_j=K$. Partition the path into consecutive
segments of these lengths, beginning at
$s_j=1+\sum_{i<j}r_i$.
For any ordered basis $N_1,\ldots,N_4$ of $V^*$, let $t_N$ be the dual
tensor whose factor on segment $j$ is
\[
 \sum_{\pi\in S_{r_j}}\sgn(\pi)
       N_{\pi(1)}\otimes\cdots\otimes N_{\pi(r_j)}.
\]
The segments are the columns of a tableau of shape $\eta$ and the initial
basis entries label its rows. By \Cref{lem:polytabloid}, their product
$t_N$ is a polytabloid, and
\[
 \mathcal T_N=\C[S_F]t_N\ \subseteq\ (V^*)^{\otimes K}
 \quad\text{is irreducible of type }\eta.
\]

First take $N_i=M_i$. Multilinearity of
\Cref{path:contraction-formula} and the definition of $t_M$ give
\begin{equation}\label{path:ordered-factorization}
 \langle u,t_M\rangle=(-1)^{m-1}[B]_{11},\qquad
 B=A_{r_1,s_1}A_{r_2,s_2}\cdots A_{r_c,s_c}.
\end{equation}
Indeed, summing independently over the permutations within each segment
is exactly the distributive expansion of the displayed ordered product.
The segments are consecutive, so no matrix factors are interchanged in
this factorization. By \Cref{path:segment}, $B$ is invertible.

The identity $\adj(X)=\tr(X)I-X$ also shows that, for every
$Q\in\GL_2(\C)$,
\[
 D_q(QXQ^{-1})=QD_q(X)Q^{-1}.
\]
Choose an eigenvector $v$ of $B$ with eigenvalue $\beta$; invertibility
gives $\beta\ne0$. Choose $Q$ taking $v$ to the first coordinate vector
and set $N_i=QM_iQ^{-1}$. Conjugation is an invertible linear operation
on $V^*$, so these $N_i$ are still an ordered basis. Each alternating
segment matrix is now $QA_{r_j,s_j}Q^{-1}$, and their ordered product is
$QBQ^{-1}$. Since $(QBQ^{-1})e_1=\beta e_1$,
\begin{equation}\label{path:nonzero-pairing}
 \langle u,t_N\rangle=(-1)^{m-1}[QBQ^{-1}]_{11}
                    =(-1)^{m-1}\beta\ne0.
\end{equation}
Only the test covectors have changed; $u$ and its endpoint letters remain
those fixed in \Cref{path:generator-sign,path:index-sum}.

The invariant evaluation pairing defines
\begin{equation}\label{path:quotient-map}
 \Phi:U_{m,2}\longrightarrow\mathcal T_N^*,\qquad
 \Phi(x)(t)=\langle x,t\rangle.
\end{equation}
For $g\in S_F$, its equivariance follows explicitly from
$\Phi(gx)(t)=\langle x,g^{-1}t\rangle=(g\Phi(x))(t)$.
By \Cref{path:nonzero-pairing} this map is nonzero on the particular
generator $u$. Its target is irreducible and, by the self-duality in
\Cref{tools:dual-sign}, has type $\eta$. Thus $\Phi$ is surjective, and
semisimplicity proves the claimed occurrence in $U_{m,2}$.
\end{proof}

\begin{example}\label{path:conjugation-example}
The change of test basis can be essential. For $m=3$ and
$\eta=(2,2,1)$, the column segments have lengths $3,2$, beginning at
positions $1,4$. Their product in \Cref{path:ordered-factorization} is
$(-6J)(2Z)=12T$, whose $(1,1)$ entry vanishes. With
\[
 Q=\begin{pmatrix}1&1\\1&2\end{pmatrix},\qquad
 QTQ^{-1}=\begin{pmatrix}1&0\\3&-1\end{pmatrix},
\]
the resulting pairing is $12$, since $(-1)^{m-1}=1$.
\end{example}

\section{Completion of the induction}
\label{sec:completion}

We now combine the two band quotients with the dominance case. The
combinatorial point is that a partition outside the dominance range either
has a long enough first row for a width-one cut, or has enough boxes in its
first four rows for a width-two cut.

\begin{samepage}
\begin{lemma}[Horizontal-strip reduction]
\label{comb:strip-reduction}
Let $m\ge2$ and $\lambda\vdash N_m$. If $\lambda$ is not dominated by
$\rho_m$, then at least one of the following holds.
\begin{enumerate}[label=\textup{(\roman*)}]
\item There is a partition $\nu\vdash N_{m-1}$ such that
      $\lambda/\nu$ is a horizontal strip of size $m$.
\item One has $m\ge5$, and there are partitions
      \[
      \nu=\lambda^{(0)}\subset\lambda^{(1)}\subset\cdots
      \subset\lambda^{(q)}=\lambda,
      \qquad q\le4,
      \]
      with $\nu\vdash N_{m-2}$ and every
      $\lambda^{(i)}/\lambda^{(i-1)}$ a nonempty horizontal strip.
      Their total size is $2m-1$.
\end{enumerate}
\end{lemma}
\end{samepage}

\begin{proof}
Unlike the deletion in \Cref{lem:young-rule}, this reduction need not
preserve a prescribed dominance inequality; a rightmost suffix suffices
to preserve the partition.
If $\lambda_1\ge m$, remove the bottom box in each of the rightmost $m$
nonempty columns. Subtracting one from this suffix of the column-height
sequence preserves weak decrease and nonnegativity. The remaining boxes
therefore form a partition $\nu$, and the deleted boxes form a horizontal
strip. Since $N_m-m=N_{m-1}$, this proves \textup{(i)}.

Suppose instead that $\lambda_1\le m-1$. Failure of domination gives an
index $k<m$ with
\begin{equation}
\label{comb:dominance-failure}
\sum_{i=1}^k\lambda_i>km-\frac{k(k-1)}2.
\end{equation}
The index is less than $m$ because the first $m$ parts of the staircase
already have sum $N_m$. For $k\le3$, however,
\[
\sum_{i=1}^k\lambda_i\le k(m-1)
\le km-\frac{k(k-1)}2.
\]
Thus $4\le k<m$, and in particular $m\ge5$. Because the parts decrease,
the mean of the first four is at least the mean of the first $k$. Hence
\begin{equation}
\label{comb:four-row-capacity}
\sum_{i=1}^4\lambda_i
\ge\frac4k\sum_{i=1}^k\lambda_i
>4m-2(k-1)\ge2m+4>2m-1.
\end{equation}

A full sweep removes the bottom box of every nonempty column. If the
current row lengths are $(\alpha_1,\alpha_2,\ldots)$, subtracting one from
all its positive column heights gives row lengths
$(\alpha_2,\alpha_3,\ldots)$. The remaining diagram is nested in the
original one, and the removed cells form a horizontal strip of size
$\alpha_1$. Four full sweeps starting from $\lambda$ would therefore
remove $\lambda_1+\cdots+\lambda_4$ boxes.

Perform full sweeps until the remaining number of boxes to be removed is
at most the number of nonempty columns, then remove that number of bottom
boxes from the rightmost columns. The suffix argument in the first
paragraph shows that this last partial sweep also leaves a partition.
By \Cref{comb:four-row-capacity}, at most four nonempty sweeps remove
exactly $2m-1$ boxes. Their remainder has size
$N_m-(2m-1)=N_{m-2}$. Reversing the sequence gives \textup{(ii)}.
\end{proof}

Four strips can be necessary. Take $m=13$ and $\lambda=(8^{11},3)$,
where the exponent denotes eleven repeated parts equal to $8$.
This partition has size $91=N_{13}$ and fails domination at $k=12$,
since its first twelve parts sum to $91>90$.
Every horizontal strip in its eight-column diagram has at most eight
boxes, so three strips cannot remove the required $25=2m-1$ boxes.
The sweeps remove $8+8+8+1=25$ boxes and leave
$\nu=(8^7,7,3)$ of size $66=N_{11}$.

\begin{proof}[Proof of \Cref{thm:cyclic-saxl}]
We induct on $m$. The module $W_1$ is the one-dimensional representation
of the one-element group, so the assertion holds at $m=1$.

Let $m\ge2$, suppose the assertion holds for all smaller positive indices,
and fix $\lambda\vdash N_m$. If $\rho_m\unrhd\lambda$, then
\Cref{prop:dominated} supplies $S^\lambda$ in $W_m$. Otherwise apply
\Cref{comb:strip-reduction}.

In case \textup{(i)}, use the band cut with $s=1$, so
$|D|=N_{m-1}$ and $|F|=m$. The width-one band module $U_{m,1}$ is
trivial. By the induction hypothesis, $S_D^\nu$ is a constituent of
$W_{m-1}$. By \Cref{lem:pieri}, $S_{B_m}^\lambda$ occurs in
\[
\Ind_{S_D\times S_F}^{S_{B_m}}
\bigl(S_D^\nu\boxtimes\mathbf1\bigr).
\]
Consequently it occurs in the quotient from \Cref{prop:band-quotient}
and hence in $W_m$.

In case \textup{(ii)}, put $K=2m-1$ and use the cut with $s=2$.
Write $b_i=|\lambda^{(i)}|-|\lambda^{(i-1)}|$, so that
$b_1+\cdots+b_q=K$ and $q\le4$. Choose disjoint subsets of $F$ of
these sizes and form the permutation module
\[
A=\Ind_{S_{b_1}\times\cdots\times S_{b_q}}^{S_F}\mathbf1
  \cong\bigoplus_{\eta\vdash K}(S_F^\eta)^{\oplus a_\eta}.
\]
Repeated application of \Cref{lem:pieri}, followed by transitivity of
induction, shows that $S_{B_m}^\lambda$ is a constituent of
$\Ind_{S_D\times S_F}^{S_{B_m}}(S_D^\nu\boxtimes A)$. Equivalently,
\begin{equation}
\label{comb:positive-multiplicity}
\begin{gathered}
\sum_{\eta\vdash K}a_\eta c_{\nu,\eta}^{\lambda}>0,\\[3pt]
c_{\nu,\eta}^{\lambda}
=\dim\Hom_{S_{B_m}}\!\left(
 S_{B_m}^\lambda,
 \Ind_{S_D\times S_F}^{S_{B_m}}
 (S_D^\nu\boxtimes S_F^\eta)\right).
\end{gathered}
\end{equation}
Here the coefficients are nonnegative integers. Moreover, every
constituent of $A$ has at most $q$ rows: starting from the empty diagram,
each of the $q$ horizontal-strip additions can introduce at most one new
row, because two new rows would both require a new box in the first
column. Thus \Cref{comb:positive-multiplicity} provides some $\eta$ with
at most four rows, $a_\eta>0$, and $c_{\nu,\eta}^{\lambda}>0$.

\begin{samepage}
By \Cref{prop:two-band}, this $S_F^\eta$ occurs in $U_{m,2}$.
By induction, $S_D^\nu$ occurs in the specified cyclic module
$W_{m-2}$. One copy of each suffices to give $S_{B_m}^\lambda$ in
\[
\Ind_{S_D\times S_F}^{S_{B_m}}(W_{m-2}\boxtimes U_{m,2}),
\]
which is a quotient of $W_m$ by \Cref{prop:band-quotient}.
Complete reducibility again gives the constituent in $W_m$ itself.
This completes the induction. Notice that case \textup{(ii)} has
$m\ge5$, so every smaller module invoked has a positive index.
\end{samepage}
\end{proof}

\begin{proof}[Proof of \Cref{thm:saxl}]
The two factors generating $W_m$ have Specht type $\rho_m$.
Thus $W_m$ is a submodule of
$S^{\rho_m}\otimes_{\C}S^{\rho_m}$ under the diagonal action.
\Cref{thm:cyclic-saxl} supplies every $S^\lambda$ in that submodule,
so $g(\rho_m,\rho_m,\lambda)>0$ for every permitted $m$ and $\lambda$.
\end{proof}

\begin{thebibliography}{HSTZ13}

\bibitem[Bes18]{bessenrodt2018}
Christine Bessenrodt,
\emph{Critical classes, Kronecker products of spin characters, and the
Saxl conjecture},
Algebraic Combinatorics \textbf{1} (2018), no.~3, 353--369.
\href{https://doi.org/10.5802/alco.18}{doi:10.5802/alco.18}.

\bibitem[BBS21]{bessenrodtbowmansutton2021}
C.~Bessenrodt, C.~Bowman, and L.~Sutton,
\emph{Kronecker positivity and 2-modular representation theory},
Transactions of the American Mathematical Society, Series B
\textbf{8} (2021), 1024--1055.
\href{https://doi.org/10.1090/btran/70}{doi:10.1090/btran/70}.

\bibitem[BWZ10]{brownwilligenburgzabrocki2010}
Andrew A.~H. Brown, Stephanie van Willigenburg, and Mike Zabrocki,
\emph{Expressions for Catalan Kronecker products},
Pacific Journal of Mathematics \textbf{248} (2010), no.~1, 31--48.
\href{https://doi.org/10.2140/pjm.2010.248.31}{doi:10.2140/pjm.2010.248.31}.

\bibitem[DG94]{drenskygiambruno1994}
Vesselin Drensky and Antonio Giambruno,
\emph{Cocharacters, codimensions and Hilbert series of the polynomial
identities for $2\times2$ matrices with involution},
Canadian Journal of Mathematics \textbf{46} (1994), no.~4, 718--733.
\href{https://doi.org/10.4153/CJM-1994-040-6}{doi:10.4153/CJM-1994-040-6}.

\bibitem[Ebr25]{ebrahimi2025}
Mahdi Ebrahimi,
\emph{The problem of deciding the positivity of Kronecker coefficients
and Saxl conjecture}, preprint, 2025.
\href{https://arxiv.org/abs/2511.03484v3}{arXiv:2511.03484v3}.

\bibitem[FH91]{fultonharris1991}
William Fulton and Joe Harris,
\emph{Representation Theory: A First Course},
Graduate Texts in Mathematics, vol.~129,
Springer-Verlag, New York, 1991.
\href{https://doi.org/10.1007/978-1-4612-0979-9}{doi:10.1007/978-1-4612-0979-9}.

\bibitem[HR23]{harmanryba2023}
Nate Harman and Christopher Ryba,
\emph{A tensor-cube version of the Saxl conjecture},
Algebraic Combinatorics \textbf{6} (2023), no.~2, 507--511.
\href{https://doi.org/10.5802/alco.267}{doi:10.5802/alco.267}.

\bibitem[HSTZ13]{heidestz2013}
Gerhard Heide, Jan Saxl, Pham Huu Tiep, and Alexandre E.~Zalesski,
\emph{Conjugacy action, induced representations and the Steinberg square
for simple groups of Lie type},
Proceedings of the London Mathematical Society (3)
\textbf{106} (2013), no.~4, 908--930.
\href{https://doi.org/10.1112/plms/pds062}{doi:10.1112/plms/pds062}.

\bibitem[Ike15]{ikenmeyer2015}
Christian Ikenmeyer,
\emph{The Saxl conjecture and the dominance order},
Discrete Mathematics \textbf{338} (2015), no.~11, 1970--1975.
\href{https://doi.org/10.1016/j.disc.2015.04.027}{doi:10.1016/j.disc.2015.04.027}.
\href{https://arxiv.org/abs/1410.6549v2}{arXiv:1410.6549v2}.

\bibitem[Jam78]{james1978}
G.~D. James,
\emph{The Representation Theory of the Symmetric Groups},
Lecture Notes in Mathematics, vol.~682,
Springer-Verlag, Berlin--Heidelberg, 1978.
\href{https://doi.org/10.1007/BFb0067708}{doi:10.1007/BFb0067708}.

\bibitem[LN25]{letelliernam2025}
Emmanuel Letellier and GyeongHyeon Nam,
\emph{The Saxl conjecture and the tensor square of unipotent characters
of $\GL_n(q)$},
Algebraic Combinatorics \textbf{8} (2025), no.~4, 1119--1140.
\href{https://doi.org/10.5802/alco.434}{doi:10.5802/alco.434}.

\bibitem[Li21]{li2021}
Xin Li,
\emph{Saxl Conjecture for triple hooks},
Discrete Mathematics \textbf{344} (2021), no.~6, article 112340.
\href{https://doi.org/10.1016/j.disc.2021.112340}{doi:10.1016/j.disc.2021.112340}.

\bibitem[LS17]{luosellke2017}
Sammy Luo and Mark Sellke,
\emph{The Saxl conjecture for fourth powers via the semigroup property},
Journal of Algebraic Combinatorics \textbf{45} (2017), no.~1, 33--80.
\href{https://arxiv.org/abs/1511.02387v2}{arXiv:1511.02387v2}.

\bibitem[Ope26]{universaltensorsquares2026}
OpenAI,
\emph{Universal Tensor Squares for Symmetric Groups},
OpenAI Math Release preprint
\href{https://github.com/openai/math/blob/main/preprints/Universal-Tensor-Squares-for-Symmetric-Groups-September-24-2026/main.pdf}{OAI:Universal-Tensor-Squares-for-Symmetric-Groups-September-24-2026},
2026.

\bibitem[PPV16]{pakpanovavallejo2016}
Igor Pak, Greta Panova, and Ernesto Vallejo,
\emph{Kronecker products, characters, \mbox{partitions}, and the tensor square
conjectures},
Advances in Mathematics \textbf{288} (2016), 702--731.
\href{https://doi.org/10.1016/j.aim.2015.11.002}{doi:10.1016/j.aim.2015.11.002}.
\href{https://arxiv.org/abs/1304.0738v2}{arXiv:1304.0738v2}.
Cited numbering follows this version.

\bibitem[Rys57]{ryser1957}
H.~J. Ryser,
\emph{Combinatorial properties of matrices of zeros and ones},
Canadian Journal of Mathematics \textbf{9} (1957), 371--377.
\href{https://doi.org/10.4153/CJM-1957-044-3}{doi:10.4153/CJM-1957-044-3}.

\end{thebibliography}

\end{document}
